% GENERATED FILE. Edit canonical lessons, then run npm run generate:books.
% canonical-source-hash: fnv1a64:b3ab26d88d63e2ed
% canonical-lessons: MR-C320-sajavne, MR-C320-utsav, MR-C320-sohla, MR-C320-yatra, MR-C320-muhurta

\chapter{Festival Days}
\label{ch:mr-festival-days}
\hlchaptermodality{320}

\begin{chapteropening}
\textbf{By the end of this chapter:} \emph{I can talk about festivals, decorations and auspicious times.}

Complete the last lesson of chapter 320: I can talk about festivals, decorations and auspicious times.
\end{chapteropening}

\section[\texorpdfstring{\mr{सजवणे} (sajavṇe)}{सजवणे (sajavṇe)}]{\mr{सजवणे} (sajavṇe) — to decorate}

\label{lesson:MR-C320-sajavne}



\begin{quote}
Before the new one: say the Marathi for a firecracker, then the Marathi for fireworks.
\end{quote}

\subsection*{You'll want to know: \mr{सजवणे}}
\textbf{\mr{सजवणे}} — \emph{sajavṇe} — \textquotedblleft{}to decorate\textquotedblright{}.

\begin{quote}
\mr{आम्ही} \mr{सणासाठी} \mr{घर} \mr{फुलांनी} \mr{सजवलं}.
\end{quote}

\begin{quote}
\emph{āmhī saṇāsāṭhī ghar phulānnī sajavlaṁ.}
\end{quote}

\begin{quote}
We decorated the house with flowers for the festival.
\end{quote}

Past, with -\mr{ने} (he …): \mr{त्याने} \mr{सजवलं}, \emph{tyāne sajavlaṁ}.

Transitive; what is decorated agrees: \mr{घर} \mr{सजवलं}, \mr{खोली} \mr{सजवली}.

\begin{cousinweb}
\textbf{\mr{सजवणे}} (\emph{sajavṇe}) is a tadbhava: CDIAL lists Marathi \emph{sajavṇe}, 'to equip', under Sanskrit \emph{sajjayati}, 'equips, prepares' (CDIAL 13091).
\end{cousinweb}

\begin{grammarlens}[title={What you've built}]
One more word for telling what happened.
\end{grammarlens}

\subsection*{Your turn}
\begin{itemize}
\raggedright
  \item \emph{Say it:} \emph{sajavṇe}
  \item \emph{Say it:} \emph{sajavṇe}, once more
  \item \emph{Say it:} \emph{āmhī saṇāsāṭhī ghar phulānnī sajavlaṁ.} — We decorated the house with flowers for the festival.
  \item \emph{Recall:} say the Marathi for a firecracker, then the Marathi for fireworks, then say \emph{sajavṇe} again
\end{itemize}

\begin{tcolorbox}[breakable,colback=teal!4,colframe=teal!35!black,title={Before you move on}]
What does \mr{सजवणे} mean? (\textquotedblleft{}To decorate\textquotedblright{}.) Say it once more.
\end{tcolorbox}
\section[\texorpdfstring{\mr{उत्सव} (utsav)}{उत्सव (utsav)}]{\mr{उत्सव} (utsav) — a festivity; a festival}

\label{lesson:MR-C320-utsav}



\begin{quote}
Before the new one: say the Marathi for fireworks, then the Marathi for to decorate.
\end{quote}

\subsection*{You'll want to know: \mr{उत्सव}}
\textbf{\mr{उत्सव}} — \emph{utsav} — \textquotedblleft{}a festivity; a festival\textquotedblright{}.

\begin{quote}
\mr{गावात} \mr{दहा} \mr{दिवस} \mr{मोठा} \mr{उत्सव} \mr{असतो}.
\end{quote}

\begin{quote}
\emph{gāvāt dahā divas moṭhā utsav asto.}
\end{quote}

\begin{quote}
There is a big festival in the village for ten days.
\end{quote}

Masculine.

\mr{सण} is a feast day; \mr{उत्सव}, the celebration around it.

\begin{cousinweb}
\textbf{\mr{उत्सव}} (\emph{utsav}) is a tatsama, taken from Sanskrit \emph{utsava}, 'a festival, a joyful occasion' (Monier-Williams).
\end{cousinweb}

\begin{grammarlens}[title={What you've built}]
One more word for telling what happened.
\end{grammarlens}

\subsection*{Your turn}
\begin{itemize}
\raggedright
  \item \emph{Say it:} \emph{utsav}
  \item \emph{Say it:} \emph{utsav}, once more
  \item \emph{Say it:} \emph{gāvāt dahā divas moṭhā utsav asto.} — There is a big festival in the village for ten days.
  \item \emph{Recall:} say the Marathi for fireworks, then the Marathi for to decorate, then say \emph{utsav} again
\end{itemize}

\begin{tcolorbox}[breakable,colback=teal!4,colframe=teal!35!black,title={Before you move on}]
What does \mr{उत्सव} mean? (\textquotedblleft{}A festivity; a festival\textquotedblright{}.) Say it once more.
\end{tcolorbox}
\section[\texorpdfstring{\mr{सोहळा} (sohḷā)}{सोहळा (sohḷā)}]{\mr{सोहळा} (sohḷā) — a grand ceremony; a celebration}

\label{lesson:MR-C320-sohla}



\begin{quote}
Before the new one: say the Marathi for to decorate, then the Marathi for a festivity.
\end{quote}

\subsection*{You'll want to know: \mr{सोहळा}}
\textbf{\mr{सोहळा}} — \emph{sohḷā} — \textquotedblleft{}a grand ceremony; a celebration\textquotedblright{}.

\begin{quote}
\mr{शाळेच्या} \mr{शताब्दीचा} \mr{सोहळा} \mr{थाटात} \mr{झाला}.
\end{quote}

\begin{quote}
\emph{śāḷecyā śatābdīcā sohḷā thāṭāt jhālā.}
\end{quote}

\begin{quote}
The school's centenary celebration was held in style.
\end{quote}

Masculine.

Plural \mr{सोहळे}.

\begin{cousinweb}
\textbf{\mr{सोहळा}} (\emph{sohḷā}) is inherited: CDIAL lists Marathi \emph{sohḷā}, 'a festive occasion', under \emph{*śubhakāla}, 'an auspicious time' (CDIAL 12533).
\end{cousinweb}

\begin{grammarlens}[title={What you've built}]
One more word for telling what happened.
\end{grammarlens}

\subsection*{Your turn}
\begin{itemize}
\raggedright
  \item \emph{Say it:} \emph{sohḷā}
  \item \emph{Say it:} \emph{sohḷā}, once more
  \item \emph{Say it:} \emph{śāḷecyā śatābdīcā sohḷā thāṭāt jhālā.} — The school's centenary celebration was held in style.
  \item \emph{Recall:} say the Marathi for to decorate, then the Marathi for a festivity, then say \emph{sohḷā} again
\end{itemize}

\begin{tcolorbox}[breakable,colback=teal!4,colframe=teal!35!black,title={Before you move on}]
What does \mr{सोहळा} mean? (\textquotedblleft{}A grand ceremony; a celebration\textquotedblright{}.) Say it once more.
\end{tcolorbox}
\section[\texorpdfstring{\mr{यात्रा} (yātrā)}{यात्रा (yātrā)}]{\mr{यात्रा} (yātrā) — a pilgrimage; a temple fair}

\label{lesson:MR-C320-yatra}



\begin{quote}
Before the new one: say the Marathi for a festivity, then the Marathi for a grand ceremony.
\end{quote}

\subsection*{You'll want to know: \mr{यात्रा}}
\textbf{\mr{यात्रा}} — \emph{yātrā} — \textquotedblleft{}a pilgrimage; a temple fair\textquotedblright{}.

\begin{quote}
\mr{दरवर्षी} \mr{हजारो} \mr{लोक} \mr{या} \mr{यात्रेला} \mr{येतात}.
\end{quote}

\begin{quote}
\emph{darvarṣī hajāro lok yā yātrelā yetāt.}
\end{quote}

\begin{quote}
Every year thousands of people come to this pilgrimage.
\end{quote}

Feminine.

\mr{यात्रेला} \mr{जाणे}, to go on a pilgrimage.

\begin{cousinweb}
\textbf{\mr{यात्रा}} (\emph{yātrā}) is a tatsama, taken from Sanskrit \emph{yātrā}, 'a going, a journey, a festive procession' (Monier-Williams).
\end{cousinweb}

\begin{grammarlens}[title={What you've built}]
One more word for telling what happened.
\end{grammarlens}

\subsection*{Your turn}
\begin{itemize}
\raggedright
  \item \emph{Say it:} \emph{yātrā}
  \item \emph{Say it:} \emph{yātrā}, once more
  \item \emph{Say it:} \emph{darvarṣī hajāro lok yā yātrelā yetāt.} — Every year thousands of people come to this pilgrimage.
  \item \emph{Recall:} say the Marathi for a festivity, then the Marathi for a grand ceremony, then say \emph{yātrā} again
\end{itemize}

\begin{tcolorbox}[breakable,colback=teal!4,colframe=teal!35!black,title={Before you move on}]
What does \mr{यात्रा} mean? (\textquotedblleft{}A pilgrimage; a temple fair\textquotedblright{}.) Say it once more.
\end{tcolorbox}
\section[\texorpdfstring{\mr{मुहूर्त} (muhūrta)}{मुहूर्त (muhūrta)}]{\mr{मुहूर्त} (muhūrta) — an auspicious time}

\label{lesson:MR-C320-muhurta}



\begin{quote}
Before the new one: say the Marathi for a grand ceremony, then the Marathi for a pilgrimage.
\end{quote}

\subsection*{You'll want to know: \mr{मुहूर्त}}
\textbf{\mr{मुहूर्त}} — \emph{muhūrta} — \textquotedblleft{}an auspicious time\textquotedblright{}.

\begin{quote}
\mr{नव्या} \mr{घरात} \mr{जाण्यासाठी} \mr{त्यांनी} \mr{चांगला} \mr{मुहूर्त} \mr{पाहिला}.
\end{quote}

\begin{quote}
\emph{navyā gharāt jāṇyāsāṭhī tyānnī cāṅglā muhūrta pāhilā.}
\end{quote}

\begin{quote}
They picked an auspicious time to move into the new house.
\end{quote}

Masculine.

\mr{मुहूर्त} \mr{पाहणे}, to fix an auspicious time.

\begin{cousinweb}
\textbf{\mr{मुहूर्त}} (\emph{muhūrta}) is a tatsama, taken from Sanskrit \emph{muhūrta}, 'a moment; a fixed span of time' (Monier-Williams).
\end{cousinweb}

\begin{grammarlens}[title={What you've built}]
One more word for telling what happened.
\end{grammarlens}

\subsection*{Your turn}
\begin{itemize}
\raggedright
  \item \emph{Say it:} \emph{muhūrta}
  \item \emph{Say it:} \emph{muhūrta}, once more
  \item \emph{Say it:} \emph{navyā gharāt jāṇyāsāṭhī tyānnī cāṅglā muhūrta pāhilā.} — They picked an auspicious time to move into the new house.
  \item \emph{Recall:} say the Marathi for a grand ceremony, then the Marathi for a pilgrimage, then say \emph{muhūrta} again
\end{itemize}

\begin{tcolorbox}[breakable,colback=teal!4,colframe=teal!35!black,title={Before you move on}]
What does \mr{मुहूर्त} mean? (\textquotedblleft{}An auspicious time\textquotedblright{}.) Say it once more.
\end{tcolorbox}
